\Needspace{8\baselineskip}
\section{Challenges and future directions}
\label{sec:future}

Six problems follow from the findings of \S\ref{sec:findings} and the evaluation gaps of
\S\ref{sec:evaluation}. Several are the medical instance of open problems the general-domain field has
stated for itself~\cite{sharkey2025open}. Each comes with an experiment that existing models and data
can support.

\textit{(1) The unmeasured loci.} Three parts of the computational path carry almost no evidence. The
first is the connector of a multimodal LLM: one study probes class labels at the encoder, the connector
and the language model in fourteen models and finds a loss of fidelity~\cite{zhu2026lost}, and none
compares the decodability or controllability of a named concept across it. The second is the set of
loci that the two-dimensional single-image path leaves out: whole-slide aggregators, volumetric
encoders that pool sub-volumes before probing, and multi-image report generators that encode time at
positions the encoder never sees. The third is the text side, where the text encoder has almost no
single-locus core claim and the linguistic prior is unmeasured after fusion. The first experiment is to
probe the same finding across RAD-DINO features, the MAIRA-2 projector output and every language-model
layer on MIMIC-CXR, and to test whether an upstream intervention reaches the report. Two cheaper
companions are to test attention-MIL weights on UNI or TITAN features against CAMELYON tiles by
occlusion, and to feed blank or shuffled images to a public chest-radiograph report generator and
report, per finding, the decodability that survives image removal.

\textit{(2) From decodable to used.} High within-study decoding accuracy is reported throughout
(\S\ref{sec:decoding}), but the control that separates a selective readout from an expressive decoder
is missing. A control task in the sense of Hewitt and Liang~\cite{hewitt2019designing} assigns each
input \emph{type} a fixed random label; language has recurring types and medical images usually do not,
so under a patient-level split a type has to cross the split (a scanner, site, protocol, view or
annotated patch class) or the control measures nothing. The selectivity of medical probes is therefore
unreported rather than poor. On the intervention side, effects are input-dependent, model-specific or
null (\S\ref{sec:intervention}), and no multimodal LLM study combines same-locus probing and
controlled steering with one behavioural end-point, although random directions cause sizeable
changes~\cite{korznikov2025the} and steering vectors are unidentifiable~\cite{venkatesh2026nonident}.
A screening diagnostic for direction-selective sensitivity would require the steering effect along a
direction $v_c$ built for concept $c$ (Eq.~(\ref{eq:effect})) to exceed the 95th percentile of the
effects of $K$ pre-specified control directions matched to $v_c$ in locus, norm and activation scale,
each averaged over held-out inputs disjoint from those that built $v_c$; the further checks that a
concept-specific contribution requires are set out in \S\ref{sec:methods}. The paired probe-and-steer
design is the natural first test, and the connector experiment of item (1) its natural setting.

\textit{(3) Reproducible discovery.} SAE features and factorised directions are unstable: absorption and
splitting are structural~\cite{chanin2024a}, and only a minority of features recur across
seeds~\cite{paulo2025sparse,sainsbury2026sparse}. A feature that does not survive retraining cannot serve
as an audit identifier, so progress needs cross-seed matching, name validation and stable identifiers.
Universal and concept-anchored dictionaries, and crosscoders across seeds or checkpoints, are the
starting points~\cite{thasarathan2025universal,nasirisarvi2025sparc,minder2025overcoming}. The experiment
is to train two differently seeded dictionaries on UNI or Virchow features, report the matched-feature
fraction, and validate the names against CAMELYON regions.

\textit{(4) Shared concept resources and evaluation standards.} Concept banks are rarely
clinician-validated (\S\ref{sec:eval-what}), and protocols vary by dataset, split, label, decoder,
width, sparsity and end-point (\S\ref{sec:eval-compare}). A shared benchmark would define
ontology-aligned vocabularies with inter-rater agreement, cover confounders, shortcuts and findings,
fix splits and decoders across public encoders and multimodal LLMs, and score faithfulness, stability
and specificity separately. Without new annotation, a first version could probe RAD-DINO and MedGemma
with VinDr-CXR and PadChest labels, UNI and CONCH with CAMELYON regions, and MONET with SkinCon. Two
reporting cards, one for decoding and one for discovery, would make results commensurable.

\textit{(5) Interpretability of adaptation.} Medical foundation models are rarely used as released, and
the evidence on what adaptation changes is thin (\S\ref{sec:geometry}); checkpoint similarity and
cross-model dictionaries are only beginning to reach medicine~\cite{minder2025overcoming,ahmed2026medseglens}.
Which directions survive site transfer, and what low-rank updates create or destroy, matters for
deployment monitoring. The experiment is to adapt a public encoder to a new site, report per-layer
similarity and probe drift, and train a crosscoder to name the changes.

\textit{(6) Prospective human and regulatory evaluation.} Medical-device guidance requires declared
inputs and outputs, performance measures, bias sources and traceable
documentation~\cite{fda2025aiguidance,euaiact2024}. Feature-space methods could document encoded
attributes for bias assessment, but no reviewed study links such a measurement to a decision, and
reader studies are rare (\S\ref{sec:eval-bench}). A trial could show or withhold a shortcut audit from
developers choosing an encoder, and record their choices and the subgroup performance of the deployed
models.

None of these experiments needs a new instrument: open suites, closed-form linear
erasure~\cite{belrose2023leace}, released dictionaries~\cite{abdulaal2024an,nooralahzadeh2026universal},
the expert-verified scorers of \S\ref{sec:eval-data}~\cite{kim2024transparent} and the models of
Table~\ref{tab:models} suffice. What is feasible depends on access: embeddings alone support calibrated probing, geometry and site measurement, and erasure before
a downstream evaluation; open encoder weights add a randomly initialised baseline and layer- and
pooling-resolved probes; and only an open encoder--connector--language-model stack supports tracing,
patching and the direction-selectivity screen, the one tier that licenses a statement about use.
